\begin{frame}{Reformas Institucionales}
\begin{itemize}
    %\item Endeudamiento elevado, déficits persistentes y presión creciente por transferencias obligatorias.
    \item Activación de la cláusula de escape (2025) revela límites del actual marco fiscal.
    \item Problema estructural: obligación de incluir todas las obligaciones legales de gasto sin contar con ingresos suficientes.
    \item Práctica recurrente: sobreestimación de ingresos para cuadrar el presupuesto.
\end{itemize}
\end{frame}

\begin{frame}{Reformas Institucionales}
\begin{itemize}
    \item Oficina de Asistencia Técnica Presupuestal (OATP): fortalecer capacidad técnica del Congreso y complementar al CARF.
    \item Validación independiente y vinculante de proyecciones de ingresos (CARF) y obligaciones legales de gasto (OATP).
    \item Mecanismo automático: Si la regla fiscal impide presupuestar todas las obligaciones legales de gasto,  declarar emergencia económica para decretar impuestos.
    \item Modificar art. 348 C.P. para que, si el Congreso no aprueba el PGN, se reconduzca el presupuesto anterior.
\end{itemize}
\end{frame}

%\begin{frame}{Reformas Institucionales}
%\begin{itemize}
%    \item Unificar presupuestos de funcionamiento e inversión en MinHacienda.
%    \item Medir ejecución presupuestal excluyendo transferencias a fiducias/patrimonios autónomos.
%    \item Reformar “pérdidas de apropiación” para permitir traslado parcial de saldos no comprometidos.
%    \item Clasificación de Transacciones de Única Vez a cargo exclusivo del CARF.
%    \item Acceso público y auditable al detalle del cierre fiscal.
%\end{itemize}
%\end{frame}
